\section{Introduction}
\label{sec:introduction}

Visual control requires an agent to connect what it sees, what it can do, and what it wants to achieve. Generative policies learn to produce action sequences directly from visual conditions~\cite{chi2023diffusion}. Latent world models instead predict the consequences of actions and support planning through a learned measure of goal proximity~\cite{zhou2024dinowm,maes2026lewm}. Both capabilities are useful, but their combination introduces a question that is easy to overlook: if an action generator already proposes a successful behavior, will a world-model planner recognize and preserve it?

This question concerns both training and inference. During training, an action prediction objective may shape features used by a dynamics predictor. During inference, an action generator may provide plausible candidates for a planner. Neither mechanism guarantees better closed-loop control. Joint objectives can favor different aspects of the representation, and model-based selection can prefer an action whose predicted outcome is more favorable than its actual outcome.

Our completed action--dynamics experiments illustrate this distinction. With a fixed 200-episode evaluation cohort per task, adding action supervision to a parallel dynamics model improves Cube planning by $7.8\pm2.6$ percentage points across three training seeds, but reduces Push-T planning by $3.8\pm0.6$ points. Initializing planning with the actor produces a different pattern: relative to random initialization, success increases by 48 points on Cube and decreases by 16 points on Reacher for the primary checkpoints. These results make a universal positive-transfer explanation untenable under the tested conditions.

We investigate a shared model in which proposal generation and consequence prediction can be examined explicitly. Fast-LeWAM uses a common visual encoder and transformer with four modes: a goal-conditioned action flow, a causal latent dynamics predictor, a goal-conditioned latent path flow, and an inverse dynamics decoder. The first two provide direct control and action-space planning. The latter two add a path-based proposal mechanism: sample intermediate latent states between the current and goal embeddings, then decode the path into an action sequence.

The path-generation and inverse-dynamics planning recipe follows LeFlow~\cite{huang2026leflow}. Our question is whether these capabilities can coexist within the action--dynamics backbone and provide a useful alternative to its own action proposals. We use the same dynamics predictor to evaluate both proposal routes. Crucially, it predicts an outcome from the proposed actions; it does not score the generated path's goal endpoint, which is fixed by construction.

The resulting framework supports three lines of investigation. First, we specify a shared architecture with distinct information constraints for action generation, causal prediction, path generation, and decoding. Second, we compare direct execution, iterative search, and selection from action or latent-path candidates within a common checkpoint. Third, we use paired control evaluations and simulator-grounded diagnostics to distinguish the availability of successful candidates from the ability to select them.

The evidence in this draft has two scopes. Completed experiments establish task-dependent behavior in the earlier action--dynamics model. The four-mode extension is an ongoing study, for which final results and matched training controls are still required. We do not infer a performance benefit from architectural sharing alone.
